
\section{曲线到射影的延拓定理}

现在，我们要用上一节的主定理——\mbox{定理~15.2.2}——来证明一个有用而具体的结果。（事实上，我们本可以早得多就证明它——只要稍加巧思，你就可以把对\mbox{定理~15.2.2} 的援引换成\mbox{练习~7.3.M}，并且一旦认识了\trDVR，便能证明它。）

作为动机，我们来考虑 \S15.2 开头那个例子的一个变体：\eqref{eq:cubic-map-failed} 是一次想要得到\trMorphism{} $\mathbb{P}^2_k\to\mathbb{P}^3_k$ 的失败尝试。这个例子用 $\mathbb{P}^2_k$ 而非 $\mathbb{P}^1_k$ 作为源，是有一番道理的。举例来说，乍看之下，由
\[
[x^7,\; x^3y^4,\; x^4y^3+4x^5y^2,\; x^5y^2+x^7]
\]
给出的\trRationalMap{} $\mathbb{P}^1_k\dashrightarrow\mathbb{P}^3_k$ 似乎并不给出一个\trMorphism{} $\mathbb{P}^1_k\to\mathbb{P}^3_k$，因为它在 $[x,y]=[0,1]$ 处没有定义。但我们其实可以把这个映射延拓到 $[0,1]$ 之上，用的是
\begin{equation*}\tag{15.3.0.1}\label{eq:curve-proj-divided}
[x^7,\; x^3y^4,\; x^4y^3+4x^5y^2,\; x^5y^2+x^7]=[x^4,\; y^4,\; xy^3+4x^2y^2,\; x^2y^2+x^4]
\end{equation*}
——我们只消通除以 $x^3$，便能叫这个\trMorphism{}讲得通。类似地，$[x^2,x^{-2}y^4,x^{-1}y^3+4y^2,y^2+x^2]$ 看上去同样在 $[x,y]=[0,1]$ 处没有定义，但你可以看出该怎么办（以及它为何能延拓成与 \eqref{eq:curve-proj-divided} 相同的映射）。这将推广成下面的结果。

\textbf{15.3.1.\ 曲线到射影的延拓定理。}\ ——\ \emph{设 $C$ 是仿射\trBaseScheme{} $S=\Spec A$ 上\trPureDimension{} $1$ 的一个\trNoetherian{}概形，而 $p\in C$ 是它的一个正则\trClosedPoint。设 $Y$ 是一个射影 $S$-\trScheme。那么，任何\trMorphism{} $C\setminus\{p\}\to Y$（$S$-\trScheme{}之间的\trMorphism）都可延拓到整个 $C$ 上。}

为预先杜绝任何混淆：这里的 $C$ 是「绝对\trPureDimension{} $1$」，而非相对\trPureDimension{} $1$。

实际使用中，我们将在 $S=\Spec k$、而 $C$ 是一个 $k$-\trVariety{}时使用这条定理。我们假定 $S$ 仿射的唯一理由是：在知道什么是射影\trMorphism{}之前（\S17.3），我们还不会知道「射影 $S$-\trScheme」的含义。而一旦知道了这条陈述的含义，下面的证明便立即推广到一般的 $S$。

注意：这样的延拓若存在，必是唯一的：$C$ 的非既约\trLocus{}是一个\trClosedSubset（练习~6.6.N）。于是，把 $C$ 换成 $p$ 的一个既约的\trOpen\trNeighborhood{}之后，我们便可以运用「既约到分离」定理~11.3.2，即：从既约概形到分离概形的\trMorphism{}由其在某个\trDense\trOpenSubset{}上的行为所确定。另一种说法是：映入分离概形的\trMorphism{}越过某个\trEffectiveCartierDivisor{}的延拓至多只有一种方式（\mbox{练习~11.3.G}）。

下面的练习表明，这些假设都是必要的。

\textbf{15.3.A.（练习）}\quad 在以下每一种情形中，证明：\trMorphism{} $C\setminus\{p\}\to Y$ 都无法延拓为\trMorphism{} $C\to Y$。

（a）\emph{$Y$ 的射影性是必要的。}设 $C=\mathbb{A}^1_k$，$p=0$，$Y=\mathbb{A}^1_k$，而 $C\setminus\{p\}\to Y$ 由「$t\mapsto 1/t$」给出。

（b）\emph{$C$ 的一维性是必要的。}设 $C=\mathbb{A}^2_k$，$p=(0,0)$，$Y=\mathbb{P}^1_k$，而 $C\setminus\{p\}\to Y$ 由 $(x,y)\mapsto[x,y]$ 给出。

（c）\emph{$C$ 的非奇异性是必要的。}设 $C=\Spec k[x,y]/(y^2-x^3)$，$p=0$，$Y=\mathbb{P}^1_k$，而 $C\setminus\{p\}\to Y$ 由 $(x,y)\mapsto[x,y]$ 给出。

我们指出：把这条（容易的）定理与（艰难的）\trProperness{}的\trValuativeCriteria（\mbox{定理~13.7.6}）结合起来，便可以绕开（棘手但并不艰难的）\mbox{\trFundamentalTheoremOfEliminationTheory~8.4.10}，得到\trProjectiveSpace{}之\trProperness{}的一个证明（见\mbox{练习~13.7.F}）。更花哨的注记：\trProperness{}的\trValuativeCriteria{}还可以用来证明：若只要求 $Y$ 是\trProper{}的，\mbox{定理~15.3.1} 依然为真，但这需要花些心思。

\textbf{15.3.2.}\ \emph{证明的中心想法。}\quad 这个证明的中心想法可以概括为「去分母」，而下面的动机示例正说明了它。假设你有一个从 $\mathbb{A}^1\setminus\{0\}$ 到\trProjectiveSpace{}的\trMorphism，而你想把它延拓到 $\mathbb{A}^1$ 上。假设这个映射由 $t\mapsto[t^4+t^{-3},t^{-2}+4t]$ 给出。那么你当然会「去分母」，把映射换成 $t\mapsto[t^7+1,t+4t^4]$。类似地，若映射由 $t\mapsto[t^2+t^3,t^2+t^4]$ 给出，你就会除以 $t^2$，得到映射 $t\mapsto[1+t,1+t^2]$。

\textbf{15.3.3.}\ \emph{证明。}\quad 我们的计划是设法把自己挪进这样一种境地：能够应用 \S15.3.2 的想法。我们先做几个快速的化归。$C$ 的非既约\trLocus{}是闭的（\mbox{练习~6.6.N}）且不含 $p$，于是把 $C$ 换成 $p$ 的一个适当的\trOpen\trNeighborhood{}之后，便可不妨设 $C$ 既约且仿射。

接下来，我们化归到 $Y=\mathbb{P}^n_A$ 的情形。选取一个\trClosedEmbedding{} $Y\to\mathbb{P}^n_A$。若该结果对 $\mathbb{P}^n$ 成立，而我们有一个\trMorphism{} $C\to\mathbb{P}^n$ 使 $C\setminus\{p\}$ 映入 $Y$，则 $C$ 也必映入 $Y$。理由：我们可以化归到源是某个\trAffine\trOpenSubset、目标是 $\mathbb{A}^n_A\subset\mathbb{P}^n_A$（从而仿射）的情形。这时，在 $Y\cap\mathbb{A}^n_A$ 上消失的那些函数拉回到 $C$ 上，便成为在 $C$ 的诸\trIrreducibleComponent{}的诸\trGenericPoint{}处消失的函数，从而在 $C$ 上处处消失（利用 $C$ 的既约性）；也就是说，$C$ 映入 $Y$。

选取 $C$ 在 $p$ 处的\trLocalRing{}中的一个\trUniformizer{} $t\in\mathfrak{m}\setminus\mathfrak{m}^2$。由于 $t$ 是 $p$ 的某个\trNeighborhood{}上的函数，把 $C$ 换成这个\trNeighborhood{}之后，便可不妨设 $t$ 是 $C$ 上的一个函数。于是 $V(t)$ 包含 $p$（因为 $t\in\mathfrak{m}$），但 $V(t)$ 不包含 $C$ 中含 $p$ 的那条\trIrreducibleComponent（因为 $t\notin\mathfrak{m}^2$）；因此，把 $C$ 换成它在 $(C\setminus V(t))\cup\{p\}$ 中 $p$ 的一个仿射\trOpen\trNeighborhood——注意 $(C\setminus V(t))\cup\{p\}\subset C$ 是开集，因为 $C$（几乎）具有\trCofiniteTopology\footnote{译注：这正是「几乎」的含义——$C$ 的真闭子集若不含整条\trIrreducibleComponent，便只是有限个点。}——我们便可不妨设 $p$ 是函数 $t$ 唯一的\trZeros（并且，当然，$t$ 在 $p$ 处以\trMultiplicity{} $1$ 消失）。

我们有一个映射 $C\setminus\{p\}\to\mathbb{P}^n_A$，由\mbox{定理~15.2.2}，它对应于 $C\setminus\{p\}$ 上的一个\trLineBundle{} $\mathscr{L}$ 以及它的 $n+1$ 个在 $C\setminus\{p\}$ 中无公共\trZeros{}的\trSection。取 $C\setminus\{p\}$ 的一个非空\trOpen{}集 $U$，使 $\mathscr{L}\cong\OO$。于是，把 $C$ 换成 $U\cup\{p\}$ 之后，我们可以把到 $\mathbb{P}^n$ 的这个映射解读为 $n+1$ 个\emph{\trRationalFunction} $f_0,\dots,f_n$，它们在除去 $p$ 之处有定义\footnote{原书此处作（defined away away from $p$），「away」重出，当为衍文，径删。——译者注}，并且在除去 $p$ 之处没有公共\trZeros。令 $N=\min_i(\mathrm{val}_p f_i)$。则 $t^{-N}f_0,\dots,t^{-N}f_n$ 是 $n+1$ 个没有公共\trZeros{}的函数。于是它们确定一个\trMorphism{} $C\to\mathbb{P}^n_A$，如愿地延拓了 $C\setminus\{p\}\to\mathbb{P}^n_A$。\hfill$\square$

\textbf{15.3.B.（练习；有用的锻炼）}\quad 设 $X$ 是一个\trNoetherian{} $k$-\trScheme，而 $Z$ 是一个\trIrreducible{}的余维数 $1$ 的\trSubvariety，其\trGenericPoint{}是 $X$ 的一个\trRegular{}点（于是\trLocalRing{} $\OO_{X,Z}$ 是一个\trDVR）。设 $\pi\colon X\dashrightarrow Y$ 是到某个射影 $k$-\trScheme{}的一个\trRationalMap。证明：该\trRationalMap{}的\trDomainOfDefinition{}包含 $Z$ 的一个\trDense\trOpenSubset。换言之，从\trNoetherian{} $k$-\trScheme{}到射影 $k$-\trScheme{}的\trRationalMap{}可以在正则的余维数 $1$ 的集合上延拓。（我们此前已经见过这条原理在起作用——见关于\trCremona{}变换的\mbox{练习~7.5.J}。）

\textbf{15.3.4.}\ \hard\ \textbf{扩展示例：正则平面\trCurve{}的\trTangentLine{}是\trSecant（secant）的极限。}

我们现在来讨论一个扩展示例，它或许能对许多不同的事情给出洞察。在学习微积分时，我们有时被教导：平面\trCurve{}的\trTangentLine{}是\trSecant{}的极限。如今我们可以在\mbox{曲线到射影的延拓定理~15.3.1}\footnote{原书此处作（The Curve-to-Projective Extension Theorem 11.1.1），编号 11.1.1 当系 15.3.1 之误（该定理即本节定理 15.3.1），径改。——译者注} 的框架下看待这件事，$\delta$ 与 $\varepsilon$ 由此变成了容易的代数。这段讨论还预示了\trDiagonalMorphism{}在理解\trDifferential{}时的重要性。

设 $k$ 是一个\trField，你起初不妨把它取为\trAlgClosed{}的。（日后你可以思考：这一切在任意\trField{}上如何讲得通。）设 $C\subset\mathbb{A}^2_k$ 是一条平面\trCurve，由方程 $f(x,y)=0$ 给出。取定 $C$ 的一个点 $p=(x_0,y_0)$。给定 $C$ 的一个点 $q=(x_1,y_1)$，设 $x_1\ne x_0$，则联结二者的直线的斜率是 $m=(y_1-y_0)/(x_1-x_0)$，而联结二者的直线的方程是
\[
(y-y_0)=\frac{y_1-y_0}{x_1-x_0}(x-x_0).
\]

更一般地（即便 $x_1=x_0$ 也可，但仍要求 $y_1\ne y_0$），联结二者的直线具有「斜率」$[x_1-x_0,y_1-y_0]\in\mathbb{P}^1_k$，而联结二者的直线的方程是 $(x_1-x_0)(y-y_0)=(y_1-y_0)(x-x_0)$。于是我们有一个映射
\begin{equation*}\tag{15.3.4.1}\label{eq:secant-slope-map}
C\setminus p\longrightarrow\mathbb{P}^1_k
\end{equation*}
它把 $q$ 送到\trSecant $\overline{pq}$ 的「斜率」。\mbox{曲线到射影的延拓定理~15.3.1} 说：只要 $p$ 是 $C$ 的一个\trSmooth{}点，这个映射便可越过 $p$ 延拓。我们来把代数算出来。请记住，此时 $x_0,y_0\in k$，因为 $p$ 是取定的点。定义 $F(x_1,y_1):=f(x_1,y_1)-f(x_0,y_0)\in k[x_1,y_1]$。

\textbf{15.3.C.（练习）}\quad 证明：$F\in(x_1-x_0,y_1-y_0)$。

于是我们可以把 $F$ 写成 $g(x_1,y_1)(x_1-x_0)+h(x_1,y_1)(y_1-y_0)$，其中 $g,h\in k[x_1,y_1]$ 的某种选取（不必唯一）。请取定这样一个 $g$ 与 $h$。

\textbf{15.3.D.（练习）}\quad 证明：若 $p=(x_0,y_0)$ 是 $C$ 的\trSmooth{}点，且带有非竖直的\trTangentVector，则 $h(x_0,y_0)\ne 0$。

\textbf{15.3.E.（练习）}\quad 证明：若 $p=(x_0,y_0)$ 是 $C$ 的\trSmooth{}点，则 $p$ 处\trTangentLine{}的斜率是 $[-h(x_0,y_0),g(x_0,y_0)]\in\mathbb{P}^1_k$，而\trTangentLine{}的方程是
\[
g(x_0,y_0)(x-x_0)+h(x_0,y_0)(y-y_0)=0.
\]

